% frontmatter/palabras/como-usar.tex -- instrucciones para el adulto
% del cuaderno de primeras palabras.
\section*{Cómo usar este cuaderno}

Este cuaderno es el paso anterior al cuaderno de frases de
\emph{Aprendo a leer}. Está pensado para una niña o un niño que ya
conoce las letras (o casi todas) y empieza a juntarlas: lee sílabas, y
quiere pasar de las sílabas a las palabras enteras. Al terminarlo,
leerá frases cortas -- justo donde empieza el cuaderno de frases.

Hay una página para cada día laborable del año, de lunes a viernes,
también en vacaciones. Cinco o diez minutos al día son suficientes.
Cada página sigue siempre la misma estructura:

\begin{itemize}[leftmargin=6mm]
\item \textbf{Fecha}, \textbf{Leído} y \textbf{Notas}: para el adulto.
  Marcar cada día si ha leído sola, con ayuda o con dificultad no
  cuesta nada, y al cabo del año es un registro de todo lo que ha
  avanzado.
\item La \textbf{palabra del día}, en el recuadro verde: arriba,
  partida en sílabas (\emph{pe · lo · ta}); debajo, entera
  (\emph{pelota}). Primero se lee sílaba a sílaba, despacio, y luego
  la palabra entera, de corrido. Al principio hay una palabra al día;
  en invierno, dos; en primavera, tres; en verano, una frase corta.
\item Una \textbf{actividad}, que casi siempre termina dibujando. Las
  instrucciones (en letra pequeña y gris) \textbf{las lee un adulto} en
  voz alta: lo único que tiene que leer la niña o el niño es lo que
  está en letra grande.
\end{itemize}

\subsection*{Las actividades}

{\small
\begin{itemize}[leftmargin=6mm, itemsep=1pt, topsep=2pt]
\item \textbf{\lblDibuja}: dibuja lo que acabas de leer. Es la manera
  más sencilla de comprobar que la palabra se ha entendido, no solo
  descifrado.
\item \textbf{\lblCompleta}: unas pocas líneas sin forma que hay que
  convertir en un dibujo (en la pera, en el sol, en la caseta de
  Toby...).
\item \textbf{\lblTraza}: repasar con el dedo o con un lápiz una letra,
  mayúscula y minúscula -- la primera letra de la palabra del día --,
  y después escribirla sola en la línea. Las 27 letras del abecedario
  aparecen a lo largo del año.
\item \textbf{\lblEscribe}: repasar la palabra del día por dentro de
  sus letras huecas, y después escribirla sin ayuda.
\item \textbf{\lblEncuentra}: rodear la palabra del día cada vez que
  aparece entre otras que se le parecen mucho (\emph{pata, pala,
  lata...}). Obliga a mirar todas las letras, no solo la primera.
\item \textbf{\lblPalmadas}: leer la palabra dando una palmada en cada
  sílaba, y colorear un círculo por palmada.
\item \textbf{\lblAdivina}: el adulto lee una adivinanza, y la niña o
  el niño dibuja la respuesta (desde el invierno, escribe también su
  nombre debajo). La respuesta suele ser una palabra que ya ha leído
  esa semana.
\item \textbf{\lblUne}: unir con una línea cada palabra con la misma
  palabra en MAYÚSCULAS -- en verano, cada animal con su sonido, cada
  palabra con su contraria o con la que rima.
\item \textbf{\lblSiNo} (solo en verano): leer frases muy cortas y
  decidir si tienen sentido (\emph{Toby ladra} -- sí; \emph{Toby
  vuela} -- no).
\item \textbf{\lblRelee} y \textbf{\lblRepasa} (los viernes): las
  palabras de toda la semana juntas, para leerlas otra vez y marcar las
  que ya salen solas. Cada diez páginas hay un pequeño cartel de
  celebración.
\end{itemize}}

\subsection*{Cómo ayudar}

{\small
\begin{itemize}[leftmargin=6mm, itemsep=1pt, topsep=2pt]
\item Si se atasca en una palabra, déjale unos segundos y, si no sale,
  léela tú señalando las sílabas con el dedo -- y que la repita. Mejor
  eso que convertir la página en una lucha.
\item No hace falta corregir cada error. Si ha leído \emph{pato} en vez
  de \emph{pata}, basta con preguntar ``¿seguro? mira el final''.
\item Si durante una semana entera lee la palabra del día a la primera,
  no hace falta esperar: se puede leer más de una página al día. Si le
  cuesta mucho, se puede repetir la página del día anterior.
\item La actividad no tiene que terminarse el mismo día. Lo importante
  es la lectura; el dibujo puede esperar a la tarde.
\end{itemize}}
\newpage
